% =====================================================================
% 附录：证明与补充理论（与 main.tex 配合编译）
% 统计学论文附录的标准结构：
%   A. 证明（按定理顺序，全部给出，不得引用不存在的中间引理）
%   B. 附加理论结果（备用引理 auxiliary lemmas，先证引理再证定理）
%   C. 额外模拟结果 / 附加真实数据分析
%   D. AI 辅助使用披露（本项目的特色，见 D 节模板）
% =====================================================================

\section{Proofs}\label{app:proof}
% 写证明前先跑 skills/stat-theory-check：
% 1) 每个引用的引理必须在此前或附录中已证；
% 2) 每个“显然/易得”必须展开或替换为可查引理；
% 3) 极限交换、求和交换必须有控制条件。

\subsection{Proof of Theorem \ref{thm:main}}\label{app:proof-thm1}
\begin{pf}
TODO：证明正文。典型结构：
Step 1（分解）：$\norm{\betaHat-\beta_0}$ 的基本不等式分解；
Step 2（控制逼近误差）：利用损失函数性质与 Assumption \ref{asm:1}；
Step 3（控制正则项误差）：KKT 条件 + 事件 $\mathcal{E}$ 的概率控制；
Step 4（合并）：在 $\lambda \asymp \sqrt{\log p/n}$ 下得到式 \eqref{eq:rate}。
（引用主文假设时用 \cref{asm:1}，须与主文实际存在的假设编号一致。）
\end{pf}

\section{Auxiliary lemmas}\label{app:lemmas}
% 所有被引用而未在正文出现的中间结果，按引理 A.1, A.2 编号并在证明前全部证完。

\section{Additional simulation results}\label{app:simsim}
% 补充场景、敏感性分析（如 λ 的取值网格、种子敏感性）。

\section{AI-assistance disclosure}\label{app:ai-disclosure}
% —— AI 辅助使用披露（按所在期刊政策调整；默认遵循 COPE 与各出版商 AI 使用规范）——
% 模板（逐条如实填写；未使用的项删除）：
%
% 1. 文献检索：使用 [工具名] 辅助检索与筛选文献，所有引用经人工逐条核对原文。
% 2. 代码：模拟与真实数据分析代码由作者编写，AI 用于代码解释与单元测试生成，
%    所有数值结果均由作者本地运行并复核（随机种子与环境见复现包）。
% 3. 理论推导：所有定理与证明由作者提出与验证；AI 用于检查证明中引理引用的一致性，
%    不作为正确性保证。
% 4. 语言润色：使用 [工具名] 对英文表达进行润色，不改变数学内容与结论。
% 5. AI 未被用于：数据伪造、结果生成、未经核验的引用插入。
%
% 提醒：多数期刊要求 AI 不能列为作者；使用披露需符合目标期刊的最新政策。
\end{document}
